\chapter{Guide d'exécution et d'exploitation}

\section{Prérequis et variables}
Le poste doit disposer de Docker Desktop, Python 3.11 à 3.13, d'un accès à PostgreSQL Ghost et Databricks, ainsi que de Power BI Desktop. Les secrets sont fournis à l'exécution et ne sont pas reproduits dans ce document.

\begin{table}[htbp]
\centering
\caption{Variables d'environnement attendues}
\label{tab:env-vars}
\small
\begin{tabularx}{\textwidth}{>{\bfseries}p{5.2cm}X}
\toprule
Variable & Usage \\
\midrule
POSTGRES\_CONNECTION & Connexion TLS à PostgreSQL \\
GHOST\_API\_KEY & Accès Ghost lorsqu'il est requis \\
DATABRICKS\_HOST & Adresse du workspace \\
DATABRICKS\_TOKEN & Jeton d'accès \\
DATABRICKS\_JOB\_ID & Job déclenché par Airflow \\
\bottomrule
\end{tabularx}
\end{table}

\section{Commandes principales}
Depuis le dossier \tech{airflow} :

\begin{lstlisting}[language=bash,caption={Démarrage Docker}]
docker compose build
docker compose up -d
docker compose ps
\end{lstlisting}

Le générateur est ensuite lancé depuis le dossier du jeu Walmart. Le port normal est \tech{http://127.0.0.1:5050}, avec recherche automatique jusqu'à 5099. Pour une démonstration : créer un client, créer une commande, relever les identifiants, déclencher le DAG si nécessaire, contrôler Gold, puis rafraîchir Power BI.

\begin{lstlisting}[language=bash,caption={Commandes dbt de contrôle}]
dbt debug
dbt source freshness
dbt run --select silver_t
dbt test --select silver_t
dbt run --select silver_b
dbt test --select silver_b
dbt run --select gold_e
dbt snapshot
dbt run --select gold_f
dbt test --select gold_f
\end{lstlisting}

\chapter{Règles du modèle Power BI}

\section{Relations}
Les cinq relations vont de la dimension côté 1 vers \tech{fact\_orders} côté N : clients par \tech{customer\_id}, produits par \tech{product\_id}, magasins par \tech{store\_id}, employés par \tech{employee\_id} et commandes par \tech{order\_id}. Le filtre croisé est simple. Les dimensions importées doivent être limitées à leur version courante.

Si Power BI signale un doublon côté 1, il faut vérifier le filtre de snapshot au lieu de forcer une relation plusieurs-à-plusieurs. Cette dernière masquerait un défaut de modèle.

\section{Mesures de départ}
\begin{lstlisting}[language={},caption={Mesures DAX recommandées}]
Revenue = SUM ( fact_orders[sales_amount] )
Orders = DISTINCTCOUNT ( fact_orders[order_id] )
Order Lines = DISTINCTCOUNT ( fact_orders[order_item_id] )
Items = SUM ( fact_orders[quantity] )
Average Order Value = DIVIDE ( [Revenue], [Orders], 0 )
Average Selling Price = DIVIDE ( [Revenue], [Items], 0 )
\end{lstlisting}

Une page exécutive contient quatre cartes, une tendance et un classement des magasins. Les pages détaillées répondent chacune à un thème. La page employés mesure l'activité enregistrée, mais ne doit pas transformer le revenu associé en jugement automatique de productivité.

\chapter{Matrice de traçabilité}

\begin{longtable}{p{1.2cm}p{4.2cm}p{4.1cm}p{3.3cm}}
\caption{Traçabilité besoin--implémentation--test}\label{tab:traceability}\\
\toprule
ID & Besoin & Implémentation & Vérification \\
\midrule
\endfirsthead
\toprule ID & Besoin & Implémentation & Vérification \\
\midrule
\endhead
F01 & Inscription automatique & Événement client Flask & Client dans dim \\
F02 & Commande cohérente & Clés actives & Zéro orphelin \\
F03 & Employé du magasin & Filtre par store & Test intertables \\
F04 & Une à cinq lignes & Produits actifs & Comptage par ordre \\
F05 & Atomicité & Transaction & Rollback forcé \\
F06 & Changements & change\_version & Curseur Bronze \\
F07 & Rejeu sûr & merge + clé unique & Test unique \\
F08 & Historique & Snapshot timestamp & Version courante \\
F09 & Grain du fait & order\_item\_id & Unique et non nul \\
F10 & BI & Schéma en étoile & Relations 1:N \\
NF01 & Secrets & Environnement & Scan du dépôt \\
NF02 & Observabilité & Logs par composant & Inspection des runs \\
\bottomrule
\end{longtable}

\section{Checklist de démonstration}
\begin{enumerate}
  \item Vérifier les conteneurs Docker.
  \item Créer un client et une commande depuis la page web.
  \item Contrôler les relations dans PostgreSQL.
  \item Déclencher le DAG ou attendre l'intervalle horaire.
  \item Confirmer le succès des onze tâches.
  \item Exécuter les requêtes de recette Gold.
  \item Rafraîchir Power BI.
  \item Contrôler \tech{Orders}, \tech{Revenue} et \tech{Items}.
\end{enumerate}

